\section{Notice that can't be turned on the person}
\label{sec:11.6}
In the Maven run notice reduced to a class and a count published afterward. Here the person can be named, and notice can do both of its jobs, but only if it is built so that it can't be used against the person it reaches. The starting model is the adverse-action notice in consumer credit. Where a compliance names, scores, locates or classifies an identifiable person, that person is told what was done, on what data, and how to respond, with a right to contest that doesn't depend on giving an account of the loss.

In an abduction and forcible transfer pipeline that model has three hazards. A notice that says \emph{you have been located and scored} is also a push to leave, and the pipeline's own contract counts leaving as an output: ImmigrationOS promised near real-time visibility into self-deportation \autocite{Ho2025ImmigrationOS}. A contest asks the person to give the agency an address and an account. And people who complain about what a system did to them are exposed to reprisal by the operator they complain about. So notice here carries four conditions:

\begin{itemize}
\item It goes through counsel or an accredited representative where the person has one, and offers one where the person doesn't.
\item It states the stay: that contesting stops the act until the contest is decided, and what the person may do in the meantime.
\item Nothing submitted in answer to it, or in a contest, may be used to locate or select the person or anyone named in it. The contest goes to counsel and the floor court under a key the agency doesn't hold, and the person proves they are the one noticed without revealing where they are, a proof in zero knowledge; the agency learns that the notice was contested and nothing that would find the person.
\item An enforcement action against a person within a set period after they contest is presumed to be reprisal, and the agency bears the burden of showing it isn't.
\end{itemize}

Until the statute carries these conditions, a deployment that holds the floor can carry them as operator terms. A deployment that can't carry them doesn't send notice to the person directly at all. It sends it to the person's counsel, or to the organizations that represent people in removal.

Those organizations are the standing notice gives to a class. They litigated the expansion of expedited removal to a stay and lost it on appeal in 2026 \autocite{MontoyaGalvez2025Cobb,Knappenberger2026DCCir}. The published classes and counts are what they need to bring a case before the individual harm has happened.

Notice is also the learning signal, and here it runs at volume. Responses to notice, contests and re-adjudications reach the system that acted, and they teach it about the people its acts fell on. They teach it with a bias: those who respond are those who could, which in an abduction and forcible transfer pipeline means those with a lawyer, an address that holds and a language the notice was written in. Sampled re-adjudication, which picks cases the affected didn't pick, corrects the bias in what the system learns. It doesn't correct the bias in what happens to people. Fourteen percent of detained people in the national study cited above had a lawyer \autocite{EaglyShafer2016Counsel}, and a statistical correction leaves the other eighty-six percent unrepresented. The remedy for the gap is the third term's counsel. Sampling is what can be done about the learning signal until the gap closes.
